% --- Archon physics formalization source begin ---
% archon:physics
% archon:covers PhyXMiniProblems/problem_phyx_mini_0329.lean
% archon:source-report reports/phyx_mini/problem_phyx_mini_0329.source.json

\chapter{Physics problem phyx\_mini\_0329}
\label{ch:PhyXMiniProblems_problem_phyx_mini_0329}

\paragraph{Problem source.}
\#\# Physical scenario

A source \$S\$ and a detector \$D\$ of radio waves are a distance \$d\$ apart on level ground. Radio waves of wavelength \$\textbackslash{}lambda\$ reach \$D\$ either along a straight path or by reflecting (bouncing) from a certain layer in the atmosphere. When the layer is at height \$H\$, the two waves reaching \$D\$ are exactly in phase. If the layer gradually rises, the phase difference between the two waves gradually shifts, until they are exactly out of phase when the layer is at height \$H + h\$.

\#\# Question

Express \$\textbackslash{}lambda\$ in terms of \$d\$, \$h\$, and \$H\$.

\#\# Answer choices

- A: A: \$2(\textbackslash{}sqrt\{2(H+h)\textasciicircum{}2 + d\textasciicircum{}2\} - \textbackslash{}sqrt\{4H\textasciicircum{}2 + d\textasciicircum{}2\})\$

- B: B: \$2(\textbackslash{}sqrt\{4(H+h)\textasciicircum{}2 + d\textasciicircum{}2\} - \textbackslash{}sqrt\{2H\textasciicircum{}2 + d\textasciicircum{}2\})\$

- C: C: \$\textbackslash{}sqrt\{4(H+h)\textasciicircum{}2 + d\textasciicircum{}2\} - \textbackslash{}sqrt\{4H\textasciicircum{}2 + d\textasciicircum{}2\}\$

- D: D: \$2(\textbackslash{}sqrt\{4(H+h)\textasciicircum{}2 + d\textasciicircum{}2\} - \textbackslash{}sqrt\{4H\textasciicircum{}2 + d\textasciicircum{}2\})\$

\#\# Figure caption (auxiliary; use the image as primary evidence)

The diagram depicts a geometric setup involving two points labeled as \textbackslash{}( S \textbackslash{}) and \textbackslash{}( D \textbackslash{}), connected by a horizontal line with an arrow indicating direction. These points are positioned at the ends of the line, which is divided into two equal parts labeled \textbackslash{}( \textbackslash{}frac\{d\}\{2\} \textbackslash{}).

Above the midpoint of the horizontal line, there are two diverging lines forming a triangle. The solid green lines and arrows illustrate a path from \textbackslash{}( S \textbackslash{}) up to the apex and then downward to \textbackslash{}( D \textbackslash{}). The apex is aligned with a vertical dashed line extending from the midpoint, perpendicular to the horizontal line.

This apex is also connected to \textbackslash{}( S \textbackslash{}) and \textbackslash{}( D \textbackslash{}) by green dashed lines, suggesting alternative paths or additional relationships. To the right, two vertical dashed lines label distances \textbackslash{}( h \textbackslash{}) and \textbackslash{}( H \textbackslash{}), representing height measurements from the horizontal line to different parts of the apex and its alignment.

\#\# Dataset metadata

Wave Properties; Physical Model Grounding Reasoning, Spatial Relation Reasoning

\paragraph{Recorded answer/context.}
D: D: \$2(\textbackslash{}sqrt\{4(H+h)\textasciicircum{}2 + d\textasciicircum{}2\} - \textbackslash{}sqrt\{4H\textasciicircum{}2 + d\textasciicircum{}2\})\$

\paragraph{Figure/image path.}
/root/proposal\_for\_physic/science-mango/phyx\_mini\_run/phyx\_data/test\_image/329.png

\paragraph{Formalization target.}
create a compiling Lean file with sorry bodies at `PhyXMiniProblems/problem\_phyx\_mini\_0329.lean`. The Lean declarations must preserve the physical quantities, dimensions or dimensional roles, figure labels, governing-law hypotheses, and final relation expressed by this problem.
Use Mathlib/Physlib names found through LeanExplore where available. If a physics API is missing, introduce faithful local abstractions rather than scalar placeholder aliases.

\begin{theorem}[Physics formalization target]
\label{thm:physics:phyx_mini_0329:target}
\lean{PhyXMiniProblems.ProblemPhyXMini0329.problem_phyx_mini_0329}
\uses{def:physics:phyx-mini-0329:phyxminiproblems-problemphyxmini0329-lengthreadout, def:physics:phyx-mini-0329:phyxminiproblems-problemphyxmini0329-atmosphericlayerinterferencesetup, def:physics:phyx-mini-0329:phyxminiproblems-problemphyxmini0329-matchesradioreflectionscenario, def:physics:phyx-mini-0329:phyxminiproblems-problemphyxmini0329-matchessuppliedfigure, def:physics:phyx-mini-0329:phyxminiproblems-problemphyxmini0329-hasphysicalparameters, def:physics:phyx-mini-0329:phyxminiproblems-problemphyxmini0329-satisfiessymmetricreflectiongeometry, def:physics:phyx-mini-0329:phyxminiproblems-problemphyxmini0329-satisfiestwopathphaselaw, def:physics:phyx-mini-0329:phyxminiproblems-problemphyxmini0329-reachesfirstoppositionafterrise, def:physics:phyx-mini-0329:phyxminiproblems-problemphyxmini0329-isuniquematchingdisplayedchoice}
The assigned autoformalize agent should translate this physics problem into Lean declarations in the covered file, with theorem and lemma proof bodies written as `by sorry`.
\end{theorem}
\begin{proof}
This is an autoformalization task, not a proof task. Produce statements that can later be proved by the physics prover without weakening the physical model.
\end{proof}

% --- Archon declaration topology begin ---
\section{Declaration topology}
\begin{definition}[Length Quantity]
  \label{def:physics:phyx-mini-0329:phyxminiproblems-problemphyxmini0329-lengthquantity}
  \lean{PhyXMiniProblems.ProblemPhyXMini0329.LengthQuantity}
  A nonnegative physical length, independent of the unit used to read it
\end{definition}
\begin{proof}
  This is the declared model definition of Length Quantity.
\end{proof}
\begin{definition}[Signed Length Quantity]
  \label{def:physics:phyx-mini-0329:phyxminiproblems-problemphyxmini0329-signedlengthquantity}
  \lean{PhyXMiniProblems.ProblemPhyXMini0329.SignedLengthQuantity}
  A signed physical length used for Cartesian coordinates in the figure
\end{definition}
\begin{proof}
  This is the declared model definition of Signed Length Quantity.
\end{proof}
\begin{definition}[length Readout]
  \label{def:physics:phyx-mini-0329:phyxminiproblems-problemphyxmini0329-lengthreadout}
  \lean{PhyXMiniProblems.ProblemPhyXMini0329.lengthReadout}
  \uses{def:physics:phyx-mini-0329:phyxminiproblems-problemphyxmini0329-lengthquantity}
  Read a nonnegative physical length in a selected length unit
\end{definition}
\begin{proof}
  This is the declared model definition of length Readout.
\end{proof}
\begin{definition}[signed Length Readout]
  \label{def:physics:phyx-mini-0329:phyxminiproblems-problemphyxmini0329-signedlengthreadout}
  \lean{PhyXMiniProblems.ProblemPhyXMini0329.signedLengthReadout}
  \uses{def:physics:phyx-mini-0329:phyxminiproblems-problemphyxmini0329-signedlengthquantity}
  Read a signed figure coordinate in a selected length unit
\end{definition}
\begin{proof}
  This is the declared model definition of signed Length Readout.
\end{proof}
\begin{definition}[Plane Point]
  \label{def:physics:phyx-mini-0329:phyxminiproblems-problemphyxmini0329-planepoint}
  \lean{PhyXMiniProblems.ProblemPhyXMini0329.PlanePoint}
  \uses{def:physics:phyx-mini-0329:phyxminiproblems-problemphyxmini0329-signedlengthquantity}
  A physical point in the vertical plane of the supplied diagram
\end{definition}
\begin{proof}
  This is the declared model definition of Plane Point.
\end{proof}
\begin{definition}[horizontal Readout]
  \label{def:physics:phyx-mini-0329:phyxminiproblems-problemphyxmini0329-horizontalreadout}
  \lean{PhyXMiniProblems.ProblemPhyXMini0329.horizontalReadout}
  \uses{def:physics:phyx-mini-0329:phyxminiproblems-problemphyxmini0329-signedlengthreadout, def:physics:phyx-mini-0329:phyxminiproblems-problemphyxmini0329-planepoint}
  Coordinate readout of a physical point in a selected length unit
\end{definition}
\begin{proof}
  This is the declared model definition of horizontal Readout.
\end{proof}
\begin{definition}[vertical Readout]
  \label{def:physics:phyx-mini-0329:phyxminiproblems-problemphyxmini0329-verticalreadout}
  \lean{PhyXMiniProblems.ProblemPhyXMini0329.verticalReadout}
  \uses{def:physics:phyx-mini-0329:phyxminiproblems-problemphyxmini0329-signedlengthreadout, def:physics:phyx-mini-0329:phyxminiproblems-problemphyxmini0329-planepoint}
  Vertical-coordinate readout of a physical point in a selected unit
\end{definition}
\begin{proof}
  This is the declared model definition of vertical Readout.
\end{proof}
\begin{definition}[Figure Point Label]
  \label{def:physics:phyx-mini-0329:phyxminiproblems-problemphyxmini0329-figurepointlabel}
  \lean{PhyXMiniProblems.ProblemPhyXMini0329.FigurePointLabel}
  Point labels visible or geometrically distinguished in the primary image
\end{definition}
\begin{proof}
  This is the declared model definition of Figure Point Label.
\end{proof}
\begin{definition}[Propagation Route]
  \label{def:physics:phyx-mini-0329:phyxminiproblems-problemphyxmini0329-propagationroute}
  \lean{PhyXMiniProblems.ProblemPhyXMini0329.PropagationRoute}
  The direct and atmospheric-reflection propagation routes
\end{definition}
\begin{proof}
  This is the declared model definition of Propagation Route.
\end{proof}
\begin{definition}[Electromagnetic Wave Kind]
  \label{def:physics:phyx-mini-0329:phyxminiproblems-problemphyxmini0329-electromagneticwavekind}
  \lean{PhyXMiniProblems.ProblemPhyXMini0329.ElectromagneticWaveKind}
  Qualitative electromagnetic-wave kind used in the problem statement
\end{definition}
\begin{proof}
  This is the declared model definition of Electromagnetic Wave Kind.
\end{proof}
\begin{definition}[Propagation Medium]
  \label{def:physics:phyx-mini-0329:phyxminiproblems-problemphyxmini0329-propagationmedium}
  \lean{PhyXMiniProblems.ProblemPhyXMini0329.PropagationMedium}
  The propagation medium between the ground apparatus and reflecting layer
\end{definition}
\begin{proof}
  This is the declared model definition of Propagation Medium.
\end{proof}
\begin{definition}[Rising Layer Figure]
  \label{def:physics:phyx-mini-0329:phyxminiproblems-problemphyxmini0329-risinglayerfigure}
  \lean{PhyXMiniProblems.ProblemPhyXMini0329.RisingLayerFigure}
  \uses{def:physics:phyx-mini-0329:phyxminiproblems-problemphyxmini0329-planepoint, def:physics:phyx-mini-0329:phyxminiproblems-problemphyxmini0329-figurepointlabel}
  Primary-image objects and marks. Coordinates are physical lengths; the Boolean fields record only what the bitmap depicts and do not impose an interference result
\end{definition}
\begin{proof}
  This is the declared model definition of Rising Layer Figure.
\end{proof}
\begin{definition}[Atmospheric Layer Interference Setup]
  \label{def:physics:phyx-mini-0329:phyxminiproblems-problemphyxmini0329-atmosphericlayerinterferencesetup}
  \lean{PhyXMiniProblems.ProblemPhyXMini0329.AtmosphericLayerInterferenceSetup}
  \uses{def:physics:phyx-mini-0329:phyxminiproblems-problemphyxmini0329-lengthquantity, def:physics:phyx-mini-0329:phyxminiproblems-problemphyxmini0329-propagationroute, def:physics:phyx-mini-0329:phyxminiproblems-problemphyxmini0329-electromagneticwavekind, def:physics:phyx-mini-0329:phyxminiproblems-problemphyxmini0329-propagationmedium, def:physics:phyx-mini-0329:phyxminiproblems-problemphyxmini0329-risinglayerfigure}
  The path lengths, wavelength, and phase observables are independent physical quantities. In particular, wavelengthLambda is not defined from the multiple-choice answer. The predicates below supply the geometric and wave laws relating these quantities
\end{definition}
\begin{proof}
  This is the declared model definition of Atmospheric Layer Interference Setup.
\end{proof}
\begin{definition}[Matches Radio Reflection Scenario]
  \label{def:physics:phyx-mini-0329:phyxminiproblems-problemphyxmini0329-matchesradioreflectionscenario}
  \lean{PhyXMiniProblems.ProblemPhyXMini0329.MatchesRadioReflectionScenario}
  \uses{def:physics:phyx-mini-0329:phyxminiproblems-problemphyxmini0329-atmosphericlayerinterferencesetup}
  Qualitative information stated in the physical scenario
\end{definition}
\begin{proof}
  This is the declared model definition of Matches Radio Reflection Scenario.
\end{proof}
\begin{definition}[Matches Supplied Figure]
  \label{def:physics:phyx-mini-0329:phyxminiproblems-problemphyxmini0329-matchessuppliedfigure}
  \lean{PhyXMiniProblems.ProblemPhyXMini0329.MatchesSuppliedFigure}
  \uses{def:physics:phyx-mini-0329:phyxminiproblems-problemphyxmini0329-lengthreadout, def:physics:phyx-mini-0329:phyxminiproblems-problemphyxmini0329-horizontalreadout, def:physics:phyx-mini-0329:phyxminiproblems-problemphyxmini0329-verticalreadout, def:physics:phyx-mini-0329:phyxminiproblems-problemphyxmini0329-atmosphericlayerinterferencesetup}
  Geometry read directly from the supplied bitmap: S, the midpoint, and D are on one horizontal ground line; both reflection points lie vertically above the midpoint; the two ground halves are each d/2; the lower apex is at H and the upper apex is an additional h higher
\end{definition}
\begin{proof}
  This is the declared model definition of Matches Supplied Figure.
\end{proof}
\begin{definition}[Has Physical Parameters]
  \label{def:physics:phyx-mini-0329:phyxminiproblems-problemphyxmini0329-hasphysicalparameters}
  \lean{PhyXMiniProblems.ProblemPhyXMini0329.HasPhysicalParameters}
  \uses{def:physics:phyx-mini-0329:phyxminiproblems-problemphyxmini0329-lengthquantity, def:physics:phyx-mini-0329:phyxminiproblems-problemphyxmini0329-lengthreadout, def:physics:phyx-mini-0329:phyxminiproblems-problemphyxmini0329-atmosphericlayerinterferencesetup}
  Positivity and nondegeneracy of all physical lengths in the scenario
\end{definition}
\begin{proof}
  This is the declared model definition of Has Physical Parameters.
\end{proof}
\begin{definition}[Satisfies Symmetric Reflection Geometry]
  \label{def:physics:phyx-mini-0329:phyxminiproblems-problemphyxmini0329-satisfiessymmetricreflectiongeometry}
  \lean{PhyXMiniProblems.ProblemPhyXMini0329.SatisfiesSymmetricReflectionGeometry}
  \uses{def:physics:phyx-mini-0329:phyxminiproblems-problemphyxmini0329-lengthquantity, def:physics:phyx-mini-0329:phyxminiproblems-problemphyxmini0329-lengthreadout, def:physics:phyx-mini-0329:phyxminiproblems-problemphyxmini0329-atmosphericlayerinterferencesetup}
  Symmetric specular-reflection geometry. At a layer height y, each leg has horizontal projection d/2 and vertical projection y, so the total reflected length is 2 * sqrt (y\textasciicircum{}2 + (d/2)\textasciicircum{}2). This general law is stated for every height and unit; it does not contain the requested wavelength formula
\end{definition}
\begin{proof}
  This is the declared model definition of Satisfies Symmetric Reflection Geometry.
\end{proof}
\begin{definition}[Satisfies Two Path Phase Law]
  \label{def:physics:phyx-mini-0329:phyxminiproblems-problemphyxmini0329-satisfiestwopathphaselaw}
  \lean{PhyXMiniProblems.ProblemPhyXMini0329.SatisfiesTwoPathPhaseLaw}
  \uses{def:physics:phyx-mini-0329:phyxminiproblems-problemphyxmini0329-lengthquantity, def:physics:phyx-mini-0329:phyxminiproblems-problemphyxmini0329-lengthreadout, def:physics:phyx-mini-0329:phyxminiproblems-problemphyxmini0329-atmosphericlayerinterferencesetup}
  Changing only the reflecting-layer height changes phase by the corresponding change of path excess divided by the wavelength. Any fixed phase contribution from reflection cancels between the two heights. Phase is measured in the dimensionless unit of wavelength cycles
\end{definition}
\begin{proof}
  This is the declared model definition of Satisfies Two Path Phase Law.
\end{proof}
\begin{definition}[Waves Are In Phase At]
  \label{def:physics:phyx-mini-0329:phyxminiproblems-problemphyxmini0329-wavesareinphaseat}
  \lean{PhyXMiniProblems.ProblemPhyXMini0329.WavesAreInPhaseAt}
  \uses{def:physics:phyx-mini-0329:phyxminiproblems-problemphyxmini0329-lengthquantity, def:physics:phyx-mini-0329:phyxminiproblems-problemphyxmini0329-atmosphericlayerinterferencesetup}
  At an in-phase height, the phase difference is an integral cycle count
\end{definition}
\begin{proof}
  This is the declared model definition of Waves Are In Phase At.
\end{proof}
\begin{definition}[Waves Are Out Of Phase At]
  \label{def:physics:phyx-mini-0329:phyxminiproblems-problemphyxmini0329-wavesareoutofphaseat}
  \lean{PhyXMiniProblems.ProblemPhyXMini0329.WavesAreOutOfPhaseAt}
  \uses{def:physics:phyx-mini-0329:phyxminiproblems-problemphyxmini0329-lengthquantity, def:physics:phyx-mini-0329:phyxminiproblems-problemphyxmini0329-atmosphericlayerinterferencesetup}
  At an out-of-phase height, the phase differs by a half-integral cycle
\end{definition}
\begin{proof}
  This is the declared model definition of Waves Are Out Of Phase At.
\end{proof}
\begin{definition}[Reaches First Opposition After Rise]
  \label{def:physics:phyx-mini-0329:phyxminiproblems-problemphyxmini0329-reachesfirstoppositionafterrise}
  \lean{PhyXMiniProblems.ProblemPhyXMini0329.ReachesFirstOppositionAfterRise}
  \uses{def:physics:phyx-mini-0329:phyxminiproblems-problemphyxmini0329-atmosphericlayerinterferencesetup, def:physics:phyx-mini-0329:phyxminiproblems-problemphyxmini0329-wavesareinphaseat, def:physics:phyx-mini-0329:phyxminiproblems-problemphyxmini0329-wavesareoutofphaseat}
  The stated observation is the first opposition reached as the layer rises: the waves start in phase at H and, after the gradual rise by h, have advanced by exactly one half-cycle and are out of phase. This is experimental phase information, not an assumption of the requested wavelength value
\end{definition}
\begin{proof}
  This is the declared model definition of Reaches First Opposition After Rise.
\end{proof}
\begin{definition}[Answer Choice]
  \label{def:physics:phyx-mini-0329:phyxminiproblems-problemphyxmini0329-answerchoice}
  \lean{PhyXMiniProblems.ProblemPhyXMini0329.AnswerChoice}
  Labels of the four wavelength formulas displayed with the problem
\end{definition}
\begin{proof}
  This is the declared model definition of Answer Choice.
\end{proof}
\begin{definition}[displayed Wavelength Readout]
  \label{def:physics:phyx-mini-0329:phyxminiproblems-problemphyxmini0329-displayedwavelengthreadout}
  \lean{PhyXMiniProblems.ProblemPhyXMini0329.displayedWavelengthReadout}
  \uses{def:physics:phyx-mini-0329:phyxminiproblems-problemphyxmini0329-lengthreadout, def:physics:phyx-mini-0329:phyxminiproblems-problemphyxmini0329-atmosphericlayerinterferencesetup, def:physics:phyx-mini-0329:phyxminiproblems-problemphyxmini0329-answerchoice}
  Scalar readout of a displayed formula in an arbitrary common length unit. The quantities H, h, and d remain dimensionful in the setup and are converted coherently before the displayed arithmetic is evaluated
\end{definition}
\begin{proof}
  This is the declared model definition of displayed Wavelength Readout.
\end{proof}
\begin{definition}[Matches Displayed Wavelength Choice]
  \label{def:physics:phyx-mini-0329:phyxminiproblems-problemphyxmini0329-matchesdisplayedwavelengthchoice}
  \lean{PhyXMiniProblems.ProblemPhyXMini0329.MatchesDisplayedWavelengthChoice}
  \uses{def:physics:phyx-mini-0329:phyxminiproblems-problemphyxmini0329-lengthreadout, def:physics:phyx-mini-0329:phyxminiproblems-problemphyxmini0329-atmosphericlayerinterferencesetup, def:physics:phyx-mini-0329:phyxminiproblems-problemphyxmini0329-answerchoice, def:physics:phyx-mini-0329:phyxminiproblems-problemphyxmini0329-displayedwavelengthreadout}
  A displayed choice agrees with the physical wavelength in every unit
\end{definition}
\begin{proof}
  This is the declared model definition of Matches Displayed Wavelength Choice.
\end{proof}
\begin{definition}[Is Unique Matching Displayed Choice]
  \label{def:physics:phyx-mini-0329:phyxminiproblems-problemphyxmini0329-isuniquematchingdisplayedchoice}
  \lean{PhyXMiniProblems.ProblemPhyXMini0329.IsUniqueMatchingDisplayedChoice}
  \uses{def:physics:phyx-mini-0329:phyxminiproblems-problemphyxmini0329-atmosphericlayerinterferencesetup, def:physics:phyx-mini-0329:phyxminiproblems-problemphyxmini0329-answerchoice, def:physics:phyx-mini-0329:phyxminiproblems-problemphyxmini0329-matchesdisplayedwavelengthchoice}
  A displayed wavelength formula is the unique matching answer
\end{definition}
\begin{proof}
  This is the declared model definition of Is Unique Matching Displayed Choice.
\end{proof}
% --- Archon declaration topology end ---
% --- Archon physics formalization source end ---
